% ===== 原文 pp.9–11（PDF pp.10–12）=====
% [原文 p.9]
\section{T-对偶}

[4, 5]

\subsection{T-对偶}

让我们考虑二维非线性 sigma 模型，其靶空间是周长为 $2\pi R$ 的圆 $S^1_R$。场变量是映射
$X:\Sigma\to\R/2\pi R\Z$，拉氏量为
\begin{equation}
  \cL=\frac{1}{2}(\partial_0X)^2-\frac{1}{2}(\partial_1X)^2. \tag{3.1}
\end{equation}
当在圆柱 $\Sigma=\R\times S^1$ 上表述、且空间坐标是周期的（$x^1\equiv x^1+2\pi$）时，存在两个守恒量。
一个是与平移对称性 $X\to X+\Delta X$ 相伴的靶动量，另一个是映射 $X$ 从空间圆到靶圆 $S^1_R$ 的绕数。
它们分别由下式度量：
\begin{equation}
  p=\frac{1}{2\pi}\int_{S^1}\partial_0X\,\dd x^1,\qquad
  w=\frac{1}{2\pi}\int_{S^1}\partial_1X\,\dd x^1. \tag{3.2}
\end{equation}
注意 $w$ 取量子化的值 $Rm$（$m$ 为某个整数）。在量子理论中，动量也量子化为 $p=l/R$（$l$ 为某个整数），
这可由如下事实看出：平移算子 $\exp(ip\Delta X)$ 在 $\Delta X=2\pi R$ 时必须是平凡的。给定 $l$ 与 $m$
的扇区中，态的能量 $E$ 与（世界面）动量 $P$ 为
\begin{equation}
  E=\frac{1}{2}\left(\left(\frac{l}{R}\right)^2+(Rm)^2\right)+E_{\mathrm{osc}},\qquad
  P=lm+P_{\mathrm{osc}}, \tag{3.3}
\end{equation}
其中 $E_{\mathrm{osc}}$ 与 $P_{\mathrm{osc}}$ 是振荡的贡献，与 $l$、$m$ 和 $R$ 无关。

我们看到，只要交换 $l$ 与 $m$，能量-动量谱就在
\begin{equation}
  R\longrightarrow\frac{1}{R} \tag{3.4}
\end{equation}
下不变。换言之，靶为 $S^1_R$ 的模型与靶为 $S^1_{1/R}$ 的模型具有相同的谱。结果表明这两个理论是等价的。
这就是\emph{T-对偶}。设 $\tX$ 是靶为 $S^1_{1/R}$ 的“对偶模型”的场，并按与 (3.2) 相同的方式
% [原文 p.10]
定义 $\tilde p$ 与 $\tilde w$。于是，原模型的动量对应于对偶的绕数，
$p\leftrightarrow\tilde w$；原模型的绕数对应于对偶的动量，$w\leftrightarrow\tilde p$。换言之，
\begin{equation}
  \partial_0X\longleftrightarrow\partial_1\tX,\qquad
  \partial_1X\longleftrightarrow\partial_0\tX. \tag{3.5}
\end{equation}

T-对偶可由如下路径积分论证导出 [6]。让我们考虑一个以映射 $X:\Sigma\to S^1_R$ 与一个形式
$J\in\Omega^1(\Sigma;\R)$ 为变量的理论，其拉氏量为
\begin{equation}
  \cL'=\frac{1}{2}\left(J_0^2-J_1^2\right)+J_1\partial_0X-J_0\partial_1X. \tag{3.6}
\end{equation}
若我们先积掉 $J$ 场——这通过对 $J$ 配方并作高斯积分完成；等价地，通过解 $J$ 的 Euler–Lagrange
方程，得到 $J_0=\partial_1X$ 与 $J_1=\partial_0X$，再把答案代回拉氏量——就得到原系统 (3.1) 中的
拉氏量 $\cL$。另一方面，若我们先积掉 $X$，则得到约束 $J_\mu=\partial_\mu\tX$，其中
$\tX:\Sigma\to S^1_{1/R}$ 是某个映射。\footnote{这里面的逻辑——特别是该 $\tX$ 具有周期性
$2\pi/R$——并不平凡，而且很有意思。它大致如下：对 $X$ 的连续部分作积分给出约束：$J$ 是闭一形式。
对 $X$ 的路径积分还包含沿 $\Sigma$ 的闭一环节对 $X$ 的绕数求和，这给出约束：$J$ 沿闭一环节的周期
是 $2\pi/R$ 的整数倍。这意味着 $J=\dd\tX$，其中 $\tX:\Sigma\to S^1_{1/R}$ 为某个映射。}
把该结果代入 $\cL'$ 中剩余的项，我们得到
\begin{equation}
  \widetilde{\cL}=\frac{1}{2}(\partial_0\tX)^2-\frac{1}{2}(\partial_1\tX)^2. \tag{3.7}
\end{equation}
这就是对偶理论的拉氏量。比较 $J$ 场的两个表达式，我们得到关系 (3.5)。这一推导的优点是：它可用来
导出更复杂情形下的 T-对偶规则，而我们将反复遇到这类情形。

\para{D-膜上的 T-对偶}

考察边界条件（即 D-膜）在 T-对偶下如何变换是富有启发的。让我们把靶为 $S^1_R$ 的 sigma 模型表述在
左半 Minkowski 时空 $\Sigma=\R\times(-\infty,0]$ 上。作为典型的边界条件，我们可以考虑 Neumann 边界
条件 (N) 与 Dirichlet 边界条件 (D)：
\begin{equation}
  \text{(N)}:\ \partial_1X|_{\partial\Sigma}=0,\qquad
  \text{(D)}:\ \partial_0X|_{\partial\Sigma}=0. \tag{3.8}
\end{equation}
由关系 (3.5) 可见，$X$ 的 Neumann 边界条件对应于 $\tX$ 的 Dirichlet 边界条件，反之亦然。确切地说，存在
% [原文 p.11]
对 (N) 与 (D) 两者而言的单参数族边界条件。(D) 意味着 $X$ 的边界值是固定的，因而由该固定位置参数化。
也就是说，(D) 由 $[x]\in S^1_R$ 参数化。我们将把此边界条件称为“$[x]$ 处的 D0-膜”。另一方面，在 (N) 之下，
我们可以向作用量添加 $S_{\mathrm{bdry}}=\int_{\partial\Sigma}a\,\dd X=\int_{\partial\Sigma}a\,\partial_0X\,\dd x^0$。
这意味着“边界粒子”——即在 $S^1_R$ 中沿轨迹 $X|_{\partial\Sigma}$ 运动的粒子——在 $U(1)$ 规范势
$A=a\,\dd x$ 下带荷。由于圆上 $U(1)$ 联络的不变量是和乐（holonomy），即 $\ee^{ia\cdot2\pi R}$，我们看到
$a$ 等价于 $a+1/R$。也就是说，(N) 由 $[2\pi a]\in S^1_{1/R}$ 参数化。我们将把此边界条件称为“缠绕在
$S^1_R$ 上、支撑着和乐为 $\ee^{ia\cdot2\pi R}$ 的 $U(1)$ 规范势的 D1-膜”。于是，T-对偶的确切陈述是：
\emph{缠绕在圆 $S^1_R$ 上的 D1-膜对应于 $S^1_{1/R}$ 中的 D0-膜，且前者的和乐对应于后者的位置。}
这可以通过考察谱来证明，或如上文那样通过路径积分论证来证明。

一个重要的教益是：对偶流形 $S^1_{1/R}$ 可以实现为原模型中缠绕在 $S^1_R$ 上的 D1-膜的模空间。
Strominger–Yau–Zaslow 的镜像对称图景正由此浮现。

\para{阿贝尔对偶}

T-对偶是所谓的“阿贝尔对偶”（Abelian duality）的一例，后者推广了四维自由 Maxwell 理论中的
电磁对偶。在每个维数 $D$，都存在一个 $k$-形式规范势 $A$ 的自由场论 $\cF^k_D(e)$，其在 $D$ 维流形 $M$
上的作用量为 $S=\frac{1}{4\pi e^2}\int_M\dd A\wedge*\dd A$。确切地说，$A$ 是一族局部定义的 $k$-形式，
使得 $\dd A$ 是整体定义的闭 $(k+1)$-形式，其上同调类取值于 $\mathrm{H}^{k+1}(M,2\pi\Z)$。阿贝尔对偶
断言：$\cF^k_D(e)$ 等价于 $\cF^{D-k-2}_D(1/e)$。这可以像 T-对偶那样用路径积分论证说明，其中 $J$ 是
一个 $(k+1)$-形式变量。T-对偶是 $\cF^0_2(1/R)\cong\cF^0_2(R)$，而电磁对偶是
$\cF^1_4(e)\cong\cF^1_4(1/e)$。

\para{推广}

我们可以考虑圆的乘积，比如 $S^1_{R_1}\times S^1_{R_2}$，作为 sigma 模型的靶空间，并对各个圆因子施加
T-对偶。例如，对第一个、第二个以及两个圆同时施加，就得到原模型与靶分别为
$S^1_{1/R_1}\times S^1_{R_2}$、$S^1_{R_1}\times S^1_{1/R_2}$ 与 $S^1_{1/R_1}\times S^1_{1/R_2}$ 的
sigma 模型之间的等价。我们也可以考虑圆 sigma 模型与任意 QFT 的乘积，并对该圆施加 T-对偶。

% [Tutor Note: p.9 式 (3.3) 可用两个守恒量直接验证：靶动量 $p=l/R$、绕数 $w=Rm$，二维玻色子的能量为
%  $E=\tfrac12 p^2+\tfrac12 w^2+E_{\mathrm{osc}}$、世界面动量为 $P=w\cdot p=lm$；把 $R\to1/R$ 且
%  $l\leftrightarrow m$ 代入，$E$ 与 $P$ 都不变，这正是正文“能谱在 (3.4) 下不变”的依据。]
% [Tutor Note: p.10 式 (3.6)→(3.1) 的配方可逐项复算：$\partial\cL'/\partial J_0=J_0-\partial_1X=0$
%  给出 $J_0=\partial_1X$，$\partial\cL'/\partial J_1=-J_1+\partial_0X=0$ 给出 $J_1=\partial_0X$；
%  代回 $\cL'$ 恰得 $\tfrac12(\partial_0X)^2-\tfrac12(\partial_1X)^2$，即 (3.1)。反向地，把
%  $J_\mu=\partial_\mu\tX$ 代入 (3.6)，交叉项 $\partial_1\tX\,\partial_0X-\partial_0\tX\,\partial_1X$
%  是全导数、积分后为零，于是剩下 (3.7)。这补出了原文“we obtain the constraint”跳过的两步。]
% [Tutor Note: p.10 式 (3.8) 与 (3.5) 的对应：由 (3.5) 的 $\partial_1X\leftrightarrow\partial_0\tX$，
%  (N) 的 $\partial_1X|_{\partial\Sigma}=0$ 等价于 $\partial_0\tX|_{\partial\Sigma}=0$，即 $\tX$ 的
%  (D)；反之 (D) 的 $\partial_0X|_{\partial\Sigma}=0$ 等价于 $\partial_1\tX|_{\partial\Sigma}=0$，
%  即 $\tX$ 的 (N)。这是正文“Neumann ↔ Dirichlet”的完整依据，也解释了为何 (N) 的模空间（和乐）与
%  (D) 的模空间（位置）在 T-对偶下互换。]
% [Tutor Note: 原文拼写笔误（译文按正确含义处理，公式不受影响）：p.10“This the the Lagrangian”（衍一
%  the）、“we obatin the relation (3.5)”（obtain）、脚注“peridicity 2π/R”（periodicity）；p.11
%  “the cricle S^1_R”（circle）、“we obtain the equaivalence of the original model”（equivalence）。]
% [Tutor Note: 和乐论证的定量核对：$U(1)$ 联络从 $a$ 变为 $a+\delta a$ 时 $\ee^{\ii a\cdot2\pi R}$
%  不变要求 $\delta a\cdot2\pi R\in2\pi\Z$，即 $a\sim a+1/R$；于是 $2\pi a$ 的周期为 $2\pi/R$，
%  取值于周长 $2\pi/R$ 的圆 $S^1_{1/R}$——这正是正文“(N) 由 $[2\pi a]\in S^1_{1/R}$ 参数化”的
%  理由，也是 (N) 的模空间恰为对偶圆的判据。]
